\documentclass[11pt]{article}
\usepackage[a4paper,margin=27mm]{geometry}
\usepackage{amsmath,amssymb,amsthm,booktabs,array,longtable}
\usepackage[T1]{fontenc}
\usepackage{lmodern}
\usepackage[colorlinks=true,linkcolor=blue,citecolor=blue,urlcolor=blue]{hyperref}
\setlength{\emergencystretch}{2em}
\newtheorem{theorem}{Theorem}
\newtheorem{lemma}[theorem]{Lemma}
\newtheorem{corollary}[theorem]{Corollary}
\newcommand{\CG}{\operatorname{CG}}
\newcommand{\E}{\operatorname{E}}
\newcommand{\diam}{\operatorname{diam}}
\newcommand{\ceil}[1]{\left\lceil #1\right\rceil}
\title{Refuting Dovgoshey--Rovenska Conjecture 4.2:\\
arXiv:2610.11596v1, Section 4, p.~13:
sharp bounds and corollaries of Petrov (2022) and Dovgoshey et al. (2026)}
\author{Haobo Ma\\\small ChronoAI Pte Ltd, Singapore
\and Wenlin Zhang\\\small National University of Singapore\\
\small ChronoAI Pte Ltd, Singapore}
\date{7 November 2026}
\begin{document}
\maketitle
\begin{abstract}
We refute Conjecture 4.2 in Section 4, p.~13, of Dovgoshey and
Rovenska, arXiv:2610.11596v1, extend their (2026) central-edge bounds,
and derive companion corollaries from Petrov's (2022) homogeneous-tree
characterization and Dovgoshey et al.'s (2026) perfect-binary-tree theorem.
A four-point ultrametric has six central edges against four in a metric.
For every $n\ge4$, the minimum number of central edges in a three-central
metric space is $2\ceil{n/2}$; within ultrametric spaces it is
$2n-4+\ceil{n/2}$.
Thus no three-central ultrametric minimizes over all metrics of the same
order. The distance-class construction also refutes Conjecture 4.1 and
answers Problem 4.3 in the same section: the smallest order of a metric
space with a $k$-element center is $1$ for $k=1$, $k$ for even $k$,
and $k+1$ for odd $k\ge3$.
Earlier counterexamples to Conjecture 4.1 occur in Love (2023) and
Bargetz et al. (2024).
More generally, at every feasible center size $k$ and order $n\ge3$,
the metric edge minimum is $(k-1)\ceil{n/2}$.
The tree corollaries give the isometry groups and sphere counts of
minimum-order ultrametrics, with explicit coverage of the related questions.
\end{abstract}
\noindent\small Licensed under Creative Commons Attribution 4.0 International
(\href{https://creativecommons.org/licenses/by/4.0/}{CC BY 4.0}).\normalsize

\section{The conjecture and a four-point counterexample}
Conjecture 4.2 of Dovgoshey and Rovenska~\cite[Section 4, p.~13]{DR}
compares a prescribed ultrametric central graph with all metric competitors.
The proposed equivalence is false: a four-point example gives six central
edges against four, and the sharp minima in the two classes differ at
every admissible order.

The all-metric statements differ from the earlier ultrametric minimum-order
theorem of Dovgoshey and Rovenska~\cite[Corollary 3.12 and Example 3.13]{UM}
and the rigidity theorem of Dovgoshey, Jaiswal, Kainth, Kaur and
Rovenska~\cite[Theorems 4.1--4.2]{BT}. The former gives minimum order
$2^{k-1}$ for a $k$-element ultrametric center; the latter characterizes
the equality spaces by perfect binary trees with one distance per level.
Section~\ref{sec:trees} states the short consequences of that prior theorem
separately from the metric edge bounds.
The homogeneous-tree characterization is already present in
Petrov~\cite[Theorem 2.2, pp.~6--7]{P}; the corresponding assertion
below is a short reproof, not a new characterization.

The center of distances, introduced by Bielas, Plewik and
Walczy\'nska~\cite{BPW}, is
\[
 C(X)=\{t\ge0:\ \text{for every }x\in X\text{ there is }y\in X
                   \text{ with }d(x,y)=t\}.
\]
In particular, $0\in C(X)$. We use the definitions of
Dovgoshey and Rovenska~\cite[Definitions 1.2 and 2.14]{DR}:
the central graph $\CG_X$ has vertex set $X$ and an edge $\{x,y\}$,
with $x\ne y$, when $d(x,y)\in C(X)$; the diametrical graph $G_X$
has an edge when $d(x,y)=\diam X$. A graph is $k$-central when
$|C(X)|=k$. Edges are unordered and loops are excluded.

Conjecture 4.2 of~\cite[Section 4, p.~13]{DR} asserts that for a
three-central finite ultrametric space $X$, $|X|\ge4$, the following
conditions are equivalent:
\begin{enumerate}
\item[(i)] $G_X$ is complete bipartite on parts $X_1,X_2$ with
$\min(|X_1|,|X_2|)=2$, and $\CG_X=G_X\cup S_1\cup S_2$, where
$S_i$ is a star on $X_i$.
\item[(ii)] $|\E(\CG_X)|\le |\E(\CG_Y)|$ for every three-central
metric space $Y$ with $|Y|=|X|$.
\end{enumerate}
The comparator in (ii) is an arbitrary metric space, not an ultrametric
space. The distinction already matters at order four.

\begin{theorem}\label{thm:four}
Conjecture 4.2 is false. Its implication $\mathrm{(i)}\Rightarrow
\mathrm{(ii)}$ fails for a four-point space.
\end{theorem}
\begin{proof}
Take $X=\{0,1,2,3\}$ with distance matrix
\[
 D_X=\begin{pmatrix}
 0&1&1/2&1\\1&0&1&1/2\\1/2&1&0&1\\1&1/2&1&0
 \end{pmatrix}.
\]
The two parts $\{0,2\}$ and $\{1,3\}$ have internal distance $1/2$
and cross-distance $1$, so this is an ultrametric. Every row has
distance set $\{0,1/2,1\}$, hence $C(X)=\{0,1/2,1\}$ and
$\CG_X=K_4$. Its diameter graph is $K_{2,2}$, and adjoining the edge
within each part gives the two required stars. Thus (i) holds and
$|\E(\CG_X)|=6$.

Let $Y=\{0,4,5,9\}\subset\mathbb R$ with its usual distance.
Its four row-distance sets are
\[
 \{0,4,5,9\},\quad\{0,1,4,5\},\quad
 \{0,1,4,5\},\quad\{0,4,5,9\}.
\]
Their intersection is $C(Y)=\{0,4,5\}$. Its central edges are precisely
$\{0,4\},\{0,5\},\{4,9\},\{5,9\}$, so
$|\E(\CG_Y)|=4<6$. Both centers have three elements, as required.
\end{proof}

\section{The sharp minimum over all metrics}
For a positive distance $t$ put
\[
 F_t=\{\{x,y\}:x\ne y,\ d(x,y)=t\}.
\]
Then $t\in C(X)$ exactly when $F_t$ is a spanning edge cover: each
vertex belongs to at least one edge. Distinct distance classes have
disjoint edge sets. This observation refines the nonempty-central-graph
bound in~\cite[Theorem 3.10]{DR} by fixing the number of central
distances.

\begin{lemma}\label{lem:cover}
If a finite metric space has $n$ points and $|C(X)|=k$, then
\[
 |\E(\CG_X)|\ge (k-1)\ceil{n/2}.
\]
\end{lemma}
\begin{proof}
Each of the $k-1$ positive central distances has a spanning edge cover
$F_t$. Counting incidences gives $2|F_t|\ge n$. Sum the resulting
$|F_t|\ge\ceil{n/2}$ over the disjoint classes.
\end{proof}

The following elementary realization uses the same narrow-band metric
principle as~\cite[proof of Theorem 3.8, formula (19)]{DR}.
Suppose $F_1,\ldots,F_s$ are edge-disjoint spanning edge covers of
$K_n$, $n\ge3$. Give all edges of $F_i$ the same length $a_i$, with
the $a_i$ distinct in $[1,2)$, and give every remaining edge its own
distinct length in $[1,2)$, different from all $a_i$. Set diagonal
distances to zero. This is a metric: in a triangle with three distinct
vertices, the longest side is less than $2$ and the other two sum to
at least $2$. Each $a_i$ is central. Each remaining length occurs on
one edge and misses a third vertex, so it is not central. Consequently
\begin{equation}\label{eq:realization}
 C(X)=\{0,a_1,\ldots,a_s\},\qquad
 \E(\CG_X)=F_1\cup\cdots\cup F_s.
\end{equation}

\begin{theorem}\label{thm:metric}
For every $n\ge4$, the minimum number of central edges among
three-central $n$-point metric spaces is $2\ceil{n/2}$.
\end{theorem}
\begin{proof}
Lemma~\ref{lem:cover} supplies the lower bound. To attain it, label
the vertices $0,\ldots,n-1$ and write the edges of their cycle as
$e_j=\{j,j+1\pmod n\}$, $0\le j<n$.
For even $n$, the even-indexed and odd-indexed cycle edges are two
disjoint perfect matchings, each of size $n/2$.
For odd $n\ge5$, let $F_1$ consist of the even-indexed edges and
let $F_2$ consist of the odd-indexed edges together with $\{0,2\}$.
Here $F_1$ covers every vertex, whereas the odd-indexed edges miss
only vertex $0$. The added chord covers it and is not a cycle edge.
Thus $F_1,F_2$ are disjoint spanning edge covers, each of size
$(n+1)/2$. Apply~\eqref{eq:realization} with $s=2$.
\end{proof}

Condition (i) in Conjecture 4.2 prescribes exactly
\begin{equation}\label{eq:old}
 2(n-2)+(2-1)+(n-2-1)=3n-6
\end{equation}
central edges. Such ultrametrics exist for every $n\ge4$: use two
parts $A,B$ of sizes $2,n-2$, give all cross-pairs distance $2$,
give pairs incident to a designated vertex in either part distance $1$,
and give pairs of distinct remaining vertices in $B$ distance $1/2$.
Every triangle within $B$ has either three sides $1/2$ or two sides
$1$; triangles meeting both parts have two sides $2$. Thus the metric
is ultrametric. Its center is $\{0,1,2\}$, since the rows in $A$
contain no other values, and every row contains these three.
Its central graph has the prescribed stars and cross-edges.
Since $3n-6>2\ceil{n/2}$ for every $n\ge4$, the failure of
Conjecture 4.2 occurs at every admissible order.

\section{Comparison restricted to ultrametrics}
A different comparison, not asserted in the source, restricts the
competitors in (ii) to ultrametrics. The prescribed $3n-6$ edges
still exceed the sharp minimum when $n\ge6$.

\begin{theorem}\label{thm:ultra}
For every $n\ge4$, the minimum number of central edges among
three-central $n$-point ultrametric spaces is
\[
 2n-4+\ceil{n/2}.
\]
Equality holds exactly when the diametrical graph is complete
bipartite with part sizes $2,n-2$, and the other positive central
distance has exactly $\ceil{n/2}$ edges.
\end{theorem}
\begin{proof}
Write $C(X)=\{0,a,D\}$, where $D=\diam X$ and $0<a<D$.
Indeed, the relation $d(x,y)<D$ partitions an ultrametric into
classes $X_1,\ldots,X_r$, $r\ge2$, and every cross-distance is $D$.
This proves directly that $D$ is central, in agreement with
\cite[Theorem 2.15]{DR}.
Every vertex has an $a$-neighbor in its own class, so $|X_i|\ge2$.
The number of diameter edges is $\sum_{i<j}|X_i||X_j|$.
Merging classes strictly decreases this sum until there are two
classes, of sizes $p,n-p$ with $2\le p\le n-2$. Also
\[
 p(n-p)-2(n-2)=(p-2)(n-p-2)\ge0.
\]
It follows that there are at least $2n-4$ diameter edges, with
equality exactly for two parts of sizes $2,n-2$. The class $F_a$
has at least $\ceil{n/2}$ edges by incidence counting.
These two bounds give the asserted inequality and equality criterion.

For attainment choose parts $A,B$ of sizes $2,n-2$. Make $A$ a
pair at distance $1$. If $|B|$ is even, partition $B$ into pairs,
each at distance $1$. If $|B|$ is odd, partition it into one triple
and pairs. On the triple give a designated vertex distance $1$
from the other two and give those two distance $1/2$; give each
pair distance $1$. Between distinct components of $B$ use distance
$3/2$, and between $A$ and $B$ use distance $2$.
The nested partition, and the fact that the two largest distances
within the triple are equal, prove the ultrametric inequality.
Every point has a distance-$1$ neighbor and a distance-$2$ neighbor;
the rows in $A$ show that the center is exactly $\{0,1,2\}$.
There are $2(n-2)$ diameter edges and $n/2$ distance-$1$ edges for
even $n$, or $2+(n-3)/2=(n+1)/2$ for odd $n$.
\end{proof}

For $n=6$ this construction has $8+3=11$ central edges, whereas
condition (i) prescribes $12$. More generally,
$2n-4+\ceil{n/2}<3n-6$ for $n\ge6$.
Finally, $2n-4>\ceil{n/2}$ for every $n\ge4$. Comparing
Theorems~\ref{thm:metric} and~\ref{thm:ultra} therefore shows that
\emph{no} three-central ultrametric satisfies the original (ii).

\section{The companion conjecture and minimum-order question}
The same distance-class viewpoint answers the other two items in
\cite[Section 4, p.~13]{DR}. Conjecture 4.1 asserts that a metric
space with $n$ points can have a complete, $n$-central graph exactly
when $n\in\{1,2,4\}$. The correct set of orders is $1$ together
with all positive even integers.

We recall the classical cyclic one-factorization~\cite[Theorem 9.1, p.~85]{H},
and give its proof to specify the metric explicitly. For even $n=2m$, use
vertices $\{\infty\}\cup\mathbb Z_q$, $q=2m-1$, and let
\begin{equation}\label{eq:factor}
 M_r=\{\{\infty,r\}\}\ \cup\
 \bigl\{\{r+j,r-j\}:1\le j\le m-1\bigr\},\qquad r\in\mathbb Z_q.
\end{equation}
Each $M_r$ is a perfect matching, since the nonzero residues modulo
the odd number $q$ occur once as $\pm j$, $1\le j\le m-1$.
Each edge incident to $\infty$ occurs once. For a finite edge
$\{u,v\}$ its round is uniquely determined by $2r=u+v\pmod q$,
and its two endpoints then have the form $r\pm j$ for exactly one
such $j$. Hence the $M_r$ partition all edges of $K_n$.

\begin{theorem}\label{thm:companion}
For $n\ge1$, a metric space of order $n$ with a complete,
$n$-central graph exists if and only if $n=1$ or $n$ is even.
\end{theorem}
\begin{proof}
Suppose $n>1$ and $|C(X)|=n$. Each vertex must realize all
$n-1$ positive central distances at its $n-1$ other vertices.
Thus each positive distance occurs exactly once in every row, and
its edges form a perfect matching. In particular $n$ is even;
all distances are central, so the graph is automatically complete.

Conversely, for even $n$ give matching $M_r$ in~\eqref{eq:factor}
length $1+r/n$, taking representatives $0\le r\le n-2$.
All lengths lie in $[1,2)$, so they define a metric. Each of the
$n-1$ positive lengths occurs at every vertex, and there are no
other lengths. Thus the center has $n$ elements and the central
graph is complete. The singleton supplies $n=1$.
\end{proof}

In particular $n=6$, with five perfect matchings at distinct lengths,
is an explicit counterexample to Conjecture 4.1.
An earlier displayed six-point example is Bargetz et al.~\cite[Example 3.11
and Figure 3]{BBKL}: every row contains every positive distance exactly once.
Love~\cite{L} had already posed this equivalent unique-distance property
and given power-of-two examples, hence counterexamples at order eight.
These antecedents are distinct from the complete corrected order
classification in Theorem~\ref{thm:companion}.

\begin{corollary}\label{cor:orders}
Let $k\ge1$. The smallest order $v(k)$ of a metric space with
$|C(X)|=k$, as asked in Problem 4.3, is
\[
 v(k)=\begin{cases}
 1,&k=1,\\
 k,&k\text{ even},\\
 k+1,&k\ge3\text{ odd}.
 \end{cases}
\]
\end{corollary}
\begin{proof}
Since $C(X)\subseteq\{d(x,y):y\in X\}$ for every fixed $x$,
we have $k\le|X|$. A singleton attains $v(1)=1$.
For even $k$, Theorem~\ref{thm:companion} attains $v(k)=k$.
For odd $k\ge3$, that theorem excludes order $k$, so
$v(k)\ge k+1$. Put $n=k+1$, which is even and at least four.
In~\eqref{eq:factor}, assign $k-1=n-2$ matchings distinct constant
lengths in $[1,3/2)$. Assign the edges of the one remaining matching
pairwise distinct lengths in $(3/2,2)$. Every constant matching length
is central. Every remaining length occurs on one edge and misses
other vertices, so it is not central. All lengths define a metric
by the same narrow-band argument. Its center therefore has exactly
$1+(k-1)=k$ elements, attaining order $k+1$.
\end{proof}

\section{All center sizes and the equality graphs}
The packing ingredients also have classical antecedents. Harary's
one-factorization~\cite[Theorem 9.1]{H} gives the even-order covers.
Andersen's cover-index criterion~\cite[abstract]{A}, applied to a complete
odd-order graph with one edge removed, gives the required odd-order packing:
exactly two vertices have minimum degree, so the cover index equals that
degree. The total edge count forces every cover in such a decomposition
to have minimum size. The proof below supplies a direct construction.

\begin{corollary}\label{cor:allk}
For $n\ge3$, the feasible center sizes are $1\le k\le n$ when $n$
is even, and $1\le k\le n-1$ when $n$ is odd. For each feasible pair,
the minimum central-edge count is
\[
 (k-1)\ceil{n/2}.
\]
At $n=1$ only $k=1$ occurs; at $n=2$ only $k=2$ occurs.
\end{corollary}
\begin{proof}
The row bound gives $k\le n$, and Theorem~\ref{thm:companion}
excludes $k=n$ for odd $n>1$. Lemma~\ref{lem:cover} gives the edge
bound. For even $n\ge4$, take any $k-1$ matchings in
\eqref{eq:factor} and apply~\eqref{eq:realization}.

For odd $n\ge3$, delete vertex $\infty$ from a one-factorization of
$K_{n+1}$. The resulting $n$ disjoint matchings $N_r$ partition
$K_n$, and $N_r$ misses exactly vertex $r$. Reserve two matchings
$N_a,N_b$. Their union consists of a path from $a$ to $b$ and even
cycles: vertices $a,b$ have degree one and all other vertices have
degree two. Orient the path toward $b$ and each cycle cyclically.
For every $r\notin\{a,b\}$, let $f_r$ be the edge directed out of
$r$. These edges are distinct and belong to the reserved matchings.
Thus $F_r=N_r\cup\{f_r\}$ are $n-2$ disjoint spanning edge covers,
each of size $(n+1)/2$. Select any $k-1$ of them and apply
\eqref{eq:realization}. Selecting none realizes $k=1$ by distinct
edge lengths, also for even $n\ge4$. The two exceptional orders
follow directly from the definition of the center.
\end{proof}

\begin{corollary}\label{cor:graph}
Among nonempty central graphs of order $n\ge2$, the minimum is
$\ceil{n/2}$. Equality graphs are $(n/2)P_2$ for even $n$ and
$P_3\cup((n-3)/2)P_2$ for odd $n$. They are bicentral.
Among all central graphs of order $n\ge3$, the minimum is zero.
\end{corollary}
\begin{proof}
A nonempty central graph has no isolated vertex: one positive
central distance already reaches every vertex. At the minimum,
$k=2$ by Lemma~\ref{lem:cover}. For even $n$, the degree sum is $n$,
so every degree is one. For odd $n$, it is $n+1$, so exactly one
degree is two and all others are one. This gives the stated graphs.
Conversely these spanning covers are realized by
\eqref{eq:realization}; order two is immediate. Distinct edge lengths
in $[1,2)$ give the empty graph when $n\ge3$.
\end{proof}
The nonempty minimum is Dovgoshey--Rovenska's
Theorem 3.10~\cite{DR}, reproved here by degree counting.
For $n=2h+1$, the component count is $P_3\cup(h-1)P_2$;
this corrects the printed $P_3\cup hP_2$ in their Theorem 3.12,
which has $n+2$ vertices. This also specifies the nonempty equality
case in their Problem 1.3; without that restriction, the minimum is zero.

\section{Corollaries of the prior ultrametric tree theorem}\label{sec:trees}
Call a finite ultrametric space MCD when
$|X|=2^{|C(X)|-1}$. Theorem 4.1 of Dovgoshey et al.~\cite{BT}
identifies its representing tree as a perfect binary tree whose
labels are constant on levels and strictly decrease toward the leaves.
Here a representing tree is the tree of recursively nested diameter
classes, with each node labeled by its diameter. Two leaves have
distance equal to the label of their last common ancestor. This
follows directly by taking the first diameter partition separating them.
An isometry permutes these classes at every node; conversely a
label-preserving tree automorphism preserves all leaf distances.
Consequently two MCD spaces with equally many points are weakly similar:
identify their binary trees and send the ordered height labels of one
to those of the other by the unique increasing bijection.
Conversely a weak similarity sends each row-distance set to its
image under that bijection, so it preserves the cardinality of their
intersection, the center.
This proves Conjecture 4.2 of~\cite[p.~20]{UM} as a consequence of
the prior classification. Its Conjectures 4.3 and 4.4 are already
proved explicitly in~\cite[Corollaries 4.3 and 4.6]{BT}:
equal centers give isometric MCD spaces, and every finite nonnegative
set containing zero is the center and distance set of an MCD space.
Additional proofs of these prescribed-center and extremal-rigidity statements
appear in Korea Superintelligence Labs~\cite[Theorem 6.1, Proposition 7.4
and Corollary 7.5]{KSI}.

The group recursion below is the binary specialization of the classical
wreath-product description of level-regular rooted-tree automorphisms;
see Ceccherini-Silberstein et al.~\cite[Sections 2.8 and 3.1, pp.~11--12]{G}.

\begin{corollary}\label{cor:binary}
Let $X$ be MCD, $|X|=2^h$, and write
$C(X)=\{0=c_0<c_1<\cdots<c_h\}$. Then $X$ is homogeneous,
meaning that its isometry group is transitive on points, and
\[
 \operatorname{Is}(X)\cong G_h,\qquad
 G_0=\{1\},\qquad
 G_h=(G_{h-1}\times G_{h-1})\rtimes C_2,
 \qquad |G_h|=2^{2^h-1},
\]
where $C_2$ swaps the two factors. Moreover, for every $p\in X$,
\[
 |\{x:d(p,x)=c_i\}|=2^{i-1}\quad(1\le i\le h).
\]
\end{corollary}
\begin{proof}
The two root subtrees are identically labeled perfect binary trees
of height $h-1$. Every isometry either preserves or swaps them,
and every pair of subtree isometries extends. This proves the group
recursion and transitivity by induction. If $g_h=|G_h|$, then
$g_0=1$ and $g_h=2g_{h-1}^2$, giving the displayed order.
The sibling subtree at height $i$ has $2^{i-1}$ leaves, precisely
the points at distance $c_i$ from $p$.
\end{proof}
This proves both parts of Conjecture 5.1 and answers Problem 5.3
of~\cite[Section 5, pp.~26--27]{BT}, as consequences of its Theorem 4.1.

\begin{corollary}\label{cor:homogeneous}
For a finite homogeneous ultrametric space $X$,
$|\operatorname{Is}(X)|\ge2^{|X|-1}$, with equality exactly when
$X$ is MCD. A finite rooted tree is the representing tree of some
homogeneous ultrametric space exactly when it has no vertex of
out-degree one and all vertices at the same level have the same
out-degree.
\end{corollary}
\begin{proof}
Suppose $|X|>1$. Transitivity permutes the diameter classes
transitively, so there are $r\ge2$ isometric classes, each of size
$m$, with $|X|=rm$. Each class $Y$ is homogeneous: an isometry
mapping two of its points must preserve their diameter class.
All isometries of the classes extend independently, and every
permutation of the classes extends using chosen identifications.
Thus $|\operatorname{Is}(X)|=r!|\operatorname{Is}(Y)|^r$.
Induct on $|X|$, using $r!\ge2^{r-1}$, with equality only for $r=2$.
Equality throughout forces binary splitting into identical subspaces
at every depth. The distinct diameter labels along a root-to-leaf
path then give $|C(X)|=1+\log_2|X|$, so $X$ is MCD.
Corollary~\ref{cor:binary} proves the converse and includes the singleton.

For the tree assertion, a transitive isometry group is transitive
on nodes at any fixed depth: map a leaf below one node to a leaf
below the other, and then map their ancestors. Hence their
out-degrees coincide, and a diameter partition has no unary node.
Conversely the stated conditions force all leaves to have one
depth, with a common branching number at each earlier level.
Assign a strictly decreasing positive label to each internal level
and zero to the leaf level. The last-common-ancestor distance is
ultrametric, and permutations of the identical subtrees act
transitively on leaves. Its diameter partitions give the prescribed tree.
\end{proof}
The group criterion is Conjecture 5.2 of~\cite[p.~27]{BT}.
The rooted-tree characterization is Conjecture 5.6 of~\cite[pp.~27--28]{BT},
reproved from the earlier characterization in
Petrov~\cite[Theorem 2.2, pp.~6--7]{P}: forgetting the level labels
gives necessity, and choosing decreasing labels on the levels gives
sufficiency.

\begin{corollary}\label{cor:spheres}
For a finite ultrametric with center $\{0=c_0<\cdots<c_h\}$,
the sphere counts $1,1,2,\ldots,2^{h-1}$ at every point hold
exactly when $X$ is MCD with $|X|=2^h$.
\end{corollary}
\begin{proof}
The forward assertion for MCD spaces is Corollary~\ref{cor:binary}.
Conversely assume the prescribed counts and $h\ge1$. The largest
central distance $c_h$ is the diameter. In its complete multipartite
graph each vertex has $2^{h-1}$ neighbors, so all $r\ge2$ parts
have the same size $t$, with $(r-1)t=2^{h-1}$. The earlier spheres
lie within a part and contain in total $2^{h-1}$ points, including
the center point. Hence $t\ge2^{h-1}$, forcing $r=2$, $t=2^{h-1}$
and $|X|=2^h$. Thus $X$ is MCD. If $h=0$, a nonsingleton finite
ultrametric would have its positive diameter in its center, so $|X|=1$.
\end{proof}
This proves Conjecture 5.5 of~\cite[p.~27]{BT} without assuming that
all distances are central in its forward implication.

For Problem 5.4 of~\cite[p.~27]{BT}, a complete answer within
homogeneous ultrametrics is also immediate. It expresses the branching
description of Delhomm\'e et al.~\cite[Section 3.1, Lemma 1;
Section 3.2, Proposition 1(3)]{D} in sphere counts. Given positive integers
$b_0,\ldots,b_h$, put $S_i=\sum_{j=0}^i b_j$.
Such a homogeneous space has these sphere counts exactly when
$b_0=1$ and $S_{i-1}$ divides $S_i$ for every $i\ge1$.
Indeed, its row-distance sets coincide, so $D(X)=C(X)$; balls of
radius $c_i$ have size $S_i$ and partition into equally sized
radius-$c_{i-1}$ balls. Conversely choose branching numbers
$S_i/S_{i-1}\ge2$ from the leaves upward and assign the increasing
heights $c_i$. The spheres have size $S_i-S_{i-1}=b_i$.
Problem 5.4 for arbitrary ultrametrics, which may have noncentral
distances, is not covered by this divisibility criterion.

\section{Coverage and verification}
\small
\begin{longtable}{@{}>{\raggedright\arraybackslash}p{.34\linewidth}>{\raggedright\arraybackslash}p{.59\linewidth}@{}}
\toprule
Source item and locator & Mathematical coverage\\
\midrule
\endhead
\cite{DR}, Conjecture 4.2, Section 4, p.~13 &
Refuted and corrected: Theorem~\ref{thm:four} and the sharp metric
replacement in Theorem~\ref{thm:metric}.\\[3pt]
\cite{DR}, Conjecture 4.1, Section 4, p.~13 &
Refuted and corrected: orders $1$ and even, Theorem~\ref{thm:companion}.
Prior counterexamples: Love (2023)~\cite{L}, Bargetz et al. (2024)~\cite{BBKL}.\\[3pt]
\cite{DR}, Problem 4.3, Section 4, p.~13 &
Answered: the all-metric minimum-order formula, Corollary~\ref{cor:orders}.\\[3pt]
\cite{DR}, Problem 1.3; Theorems 3.10, 3.12, pp.~1, 8, 11 &
Nonempty edge bound reproved; prior proof Dovgoshey--Rovenska (2026).
Odd equality component count corrected in Corollary~\ref{cor:graph}.
The unrestricted minimum is zero for $n\ge3$.\\[3pt]
\cite{DR}, Theorem 3.8, pp.~6--8; Theorem 3.13, p.~12 &
Prior proofs cited: narrow-band realization and the ultrametric
$n-1$ bound. Fixed-center extensions are
Theorem~\ref{thm:ultra} and Corollary~\ref{cor:allk}.\\[3pt]
Ultrametric-only variant of Conjecture 4.2(ii) &
Not asserted in the source; the proposed star minimum fails at $n=6$.
Theorem~\ref{thm:ultra} proves the corrected extremum.\\[3pt]
\cite{UM}, Corollary 3.12 and Example 3.13, p.~19 &
Prior proof cited: Dovgoshey--Rovenska (26 March 2026), ultrametric
minimum order $2^{k-1}$.\\[3pt]
\cite{UM}, Conjectures 4.2--4.4, p.~20 &
Prior proof cited: Dovgoshey et al. (10 September 2026),
Theorem 4.1 and Corollaries 4.3 and 4.6 of~\cite{BT}.
The weak-similarity consequence is written out in Section~\ref{sec:trees}.\\[3pt]
\cite{KSI}, Theorem 6.1 and Proposition 7.4--Corollary 7.5, pp.~7--8 &
Additional proofs cited: prescribed centers, extremal rigidity,
and forward sphere counts (Theorem 7.3).\\[3pt]
\cite{BT}, Conjecture 5.1, pp.~26--27 &
Proved as tree-theorem corollaries: homogeneity and group order,
Corollary~\ref{cor:binary}.\\[3pt]
\cite{BT}, Problem 5.3, p.~27 &
Answered: the isometry-group recursion, Corollary~\ref{cor:binary},
specializing the classical tree-group formula~\cite{G}.\\[3pt]
\cite{BT}, Conjecture 5.2, p.~27 &
Proved: the isometry-group criterion, Corollary~\ref{cor:homogeneous}.\\[3pt]
\cite{BT}, Conjecture 5.6, pp.~27--28 &
Reproved; prior proof Petrov (2022); arXiv version (2024), Theorem 2.2.
Corollary~\ref{cor:homogeneous} states its unlabeled-tree consequence.\\[3pt]
\cite{BT}, Conjecture 5.5, p.~27 &
Proved: sphere-count characterization, Corollary~\ref{cor:spheres}.
Forward counts also follow from~\cite[Theorem 7.3]{KSI}.\\[3pt]
\cite{BT}, Problem 5.4, p.~27 &
Not covered in full. The homogeneous subcase has the divisibility
criterion proved in Section~\ref{sec:trees}.\\
\bottomrule
\end{longtable}
\normalsize

The pinned Lean file
\href{https://github.com/the-omega-institute/trureturing/blob/8aa554ef66d0c1970b81b49632cd8a73a09ee34f/D5/S3/Combinatorics/Geometry/DovgosheyRovenskaCentralEdges.lean}{\texttt{DovgosheyRovenskaCentralEdges.lean}}
contains the declaration
\begin{center}\small
\path{D5.S3.Combinatorics.Geometry.DovgosheyRovenskaCentralEdges.result}
\end{center}
with statement $\neg\texttt{claim}$. Here \texttt{claim} is the
equivalence in Conjecture 4.2 with the source's metric comparator.
The declaration uses the same four-point spaces as
Theorem~\ref{thm:four}: the induced $2$-adic metric on $0,1,2,3$
and the real metric on $0,4,5,9$. Its scope is the refutation,
not the all-order bounds or the companion results.
All the additional extrema, matching constructions, equality graphs,
and tree corollaries have written proofs above. The latter use the
prior classification in~\cite[Theorem 4.1]{BT}.

The accompanying \texttt{verify.py} checks the four-point matrices,
the three all-order central-edge constructions for $4\le n\le32$,
the matching constructions for $1\le k\le24$, and all symmetric
distance tables on two to four vertices with off-diagonal values
in $\{1,2,3\}$. It also checks every feasible center size for
$3\le n\le20$, binary sphere counts through height five,
and all isometries of the binary examples through height three.
Arithmetic is exact. These finite checks support
the examples; the all-parameter conclusions follow from the proofs.

\begin{thebibliography}{99}
\bibitem{BPW}
W.~Bielas, S.~Plewik and M.~Walczy\'nska,
\emph{On the center of distances}, European Journal of Mathematics
\textbf{4} (2018), 687--698.
\href{https://doi.org/10.1007/s40879-017-0199-4}{doi:10.1007/s40879-017-0199-4}.

\bibitem{DR}
O.~Dovgoshey and O.~Rovenska,
\emph{On the Graph Corresponding to Center of Distances of Metric Space},
\href{https://arxiv.org/abs/2610.11596v1}{arXiv:2610.11596v1}
(8 October 2026).

\bibitem{UM}
O.~Dovgoshey and O.~Rovenska,
\emph{On the center of distances of finite ultrametric spaces},
\href{https://arxiv.org/abs/2603.25850v1}{arXiv:2603.25850v1}
(26 March 2026).

\bibitem{BT}
O.~Dovgoshey, N.~S.~Jaiswal, S.~P.~S.~Kainth, G.~Kaur and O.~Rovenska,
\emph{Maximal Center of Distances of Finite Ultrametric Spaces and
Perfect Binary Trees},
\href{https://arxiv.org/abs/2609.11384v1}{arXiv:2609.11384v1}
(10 September 2026).

\bibitem{P}
E.~A.~Petrov,
\emph{Hereditary properties of finite ultrametric spaces},
Journal of Mathematical Sciences \textbf{264} (2022), 423--440.
\href{https://doi.org/10.1007/s10958-022-06009-1}{doi:10.1007/s10958-022-06009-1}.
Author version:
\href{https://arxiv.org/abs/2412.17421v1}{arXiv:2412.17421v1}
(23 December 2024).

\bibitem{L}
J.~Love,
\emph{Metric spaces where each point sees each distance exactly once},
Mathematics Stack Exchange (7 February 2023), question 4634500.
\url{https://math.stackexchange.com/questions/4634500}.

\bibitem{BBKL}
C.~Bargetz, A.~Barto\v{s}, W.~Kubi\'s and F.~Luggin,
\emph{Homogeneous isosceles-free spaces},
\href{https://arxiv.org/abs/2305.03163v2}{arXiv:2305.03163v2}
(13 March 2024), Example 3.11 and Figure 3.

\bibitem{H}
F.~Harary, \emph{Graph Theory}, Addison-Wesley, Reading, MA (1969),
Chapter 9, Theorem 9.1, p.~85.

\bibitem{A}
L.~D. Andersen, \emph{Lower bounds on the cover-index of a graph},
Discrete Mathematics \textbf{25}(3) (1979), 199--210.
\href{https://doi.org/10.1016/0012-365X(79)90076-1}{doi:10.1016/0012-365X(79)90076-1}.

\bibitem{G}
T.~Ceccherini-Silberstein, D.~D'Angeli, A.~Donno, F.~Scarabotti and F.~Tolli,
\emph{Finite Gelfand pairs: examples and applications},
in \emph{Ischia Group Theory 2008}, World Scientific (2009), 7--41.
\url{https://www.researchgate.net/publication/228539253_Finite_Gelfand_pairs_Examples_and_applications}.

\bibitem{D}
C.~Delhomm\'e, C.~Laflamme, M.~Pouzet and N.~Sauer,
\emph{On homogeneous ultrametric spaces},
\href{https://arxiv.org/abs/1509.04346v2}{arXiv:1509.04346v2}
(29 November 2016; first version 14 September 2015).

\bibitem{KSI}
Korea Superintelligence Labs,
\emph{The center of distances of a finite ultrametric space:
centered spheres, prescribed centers, and extremal rigidity}, version 1.
\url{https://ideosphere.ai/papers/ultra-center-log2-v1.pdf}.
\end{thebibliography}
\end{document}
